\begin{abstract}
Exploration is a primary bottleneck in reinforcement learning when task rewards are sparse, delayed, or otherwise uninformative.
Intrinsic-reward methods based on prediction error can encourage novelty but may overemphasize irreducible stochasticity, while global progress signals can miss localized model improvement.
We introduce the \emph{GLPE family} (\textbf{G}ated \textbf{L}earning-\textbf{P}rogress \textbf{E}xploration), intrinsic rewards that target \emph{learnable} novelty by combining two signals derived from a learned latent dynamics model: region-local learning progress---estimated from changes in latent prediction error within an online partition of feature space---and feature-space impact.
\textbf{GLPE (no gate)} uses this combined score directly, yielding a simple region-aware bonus with RMS-normalized scaling.
\textbf{GLPE} augments the same score with a region-specific gate that suppresses intrinsic reward in areas that remain high-error yet exhibit little progress, serving as a targeted guardrail against unproductive curiosity.
We evaluate both variants with a common PPO backbone on five Gymnasium control benchmarks and compare against widely used intrinsic-reward baselines, reporting performance under both interaction- and wall-clock-normalized views.
Across the suite, \textbf{GLPE (no gate)} is a consistently competitive default, while gating is most beneficial in sparse-reward regimes and trades off additional intrinsic-computation overhead.
Ablations and diagnostics clarify the roles of progress, impact, and compute overhead in practice.
\end{abstract}
